%%% Copyright (C) [2026] [PWr * WIT * Wydziałowa Komisja d/s Jakości Kształcenia]
%%% Zmodyfikowany i zoptymalizowany (2026)
%%% ---------------------------------------------------------------------------

%%% KONFIGURACJA EKSTRAKCJI TEKSTU (PDF COPY-PASTE)
\input{glyphtounicode} 
\pdfgentounicode=1     % Wymagane przez JSA (Jednolity System Antyplagiatowy)

%%% 1. KODOWANIE I JĘZYK
\RequirePackage[T1]{fontenc}
%\RequirePackage[utf8]{inputenc} % Opcjonalne w nowych systemach (domyślne)
\RequirePackage[english,polish]{babel} % Kolejność eng -> pl
\RequirePackage{csquotes} % Wymagane przez biblatex

%%% 2. MATEMATYKA I CZCIONKI
%%% WIT / Zasady Grebera: Times New Roman 12 pt w całym dokumencie.
%%% newtx* = nowoczesny klon Times (tekst + matematyka w tej samej rodzinie).
%%% Wyjątek: listingi/kod — \ttfamily (maszynowa), zgodnie z praktyką techniczną.
\RequirePackage{amsmath}
\RequirePackage{mathtools}
\RequirePackage{amsthm}

\RequirePackage{newtxtext}
\RequirePackage{newtxmath}
\RequirePackage{microtype} % Poprawia optyczny wygląd tekstu
\renewcommand{\familydefault}{\rmdefault} % serif wszędzie (bez sffamily na tytule)

%%% 3. GEOMETRIA I GRAFIKA
%%% wymogi_edytorskie: 25 mm + 5 mm na oprawę (lewa / inner)
\RequirePackage{geometry}
\geometry{
  a4paper,
  top=25mm, bottom=25mm,
  inner=25mm, outer=25mm,
  bindingoffset=5mm
}
\RequirePackage{graphicx}
\RequirePackage{xcolor}
\RequirePackage{float} % Dla opcji pozycjonowania [H]

%%% 4. NARZĘDZIA POMOCNICZE
\RequirePackage{etoolbox}
\RequirePackage{indentfirst} % wcięcie także pierwszego akapitu po tytule

%%% 5. FORMATOWANIE TEKSTU I NAGŁÓWKÓW
%%% wymogi_edytorskie: rozdział 14 pt, podrozdział 13 pt, 2. rząd 12 pt; interlinia 1;
%%% Zasady Grebera: justowanie, bez kropek po tytułach rozdziałów
\RequirePackage{titlesec}
\RequirePackage{caption}
\RequirePackage{subcaption}
\RequirePackage{setspace}

\setlength{\parindent}{0.7cm}
\singlespacing % wymogi_edytorskie: interlinia pojedyncza (klasa report justuje tekst)

% Rozmiary WIT (Times / newtx): 14 / 13 / 12 pt
\newcommand{\WitChapterFont}{\normalfont\bfseries\fontsize{14}{17}\selectfont}
\newcommand{\WitSectionFont}{\normalfont\bfseries\fontsize{13}{16}\selectfont}
\newcommand{\WitSubsectionFont}{\normalfont\bfseries\fontsize{12}{14.5}\selectfont}

\titleformat{\chapter}[display]
  {\WitChapterFont}
  {\chaptertitlename\ \thechapter}{0.5ex}
  {\WitChapterFont}
\titlespacing*{\chapter}{0pt}{-20pt}{20pt}

\titleformat{\section}{\WitSectionFont}{\thesection.}{1em}{}
\titlespacing*{\section}{0pt}{2.5ex plus 1ex minus .2ex}{1.3ex plus .2ex}

\titleformat{\subsection}{\WitSubsectionFont}{\thesubsection.}{1em}{}
\titlespacing*{\subsection}{0pt}{2.25ex plus 1ex minus .2ex}{1.0ex plus .2ex}

\titleformat{\subsubsection}{\WitSubsectionFont}{\thesubsubsection.}{1em}{}
\titlespacing*{\subsubsection}{0pt}{2.0ex plus 1ex minus .2ex}{0.8ex plus .2ex}

% Numeracja ciągła wewnątrz rozdziałów (Tabela 2.1 / Rys. 2.1)
\counterwithin{figure}{chapter}
\counterwithin{table}{chapter}
\numberwithin{equation}{chapter}

% Podpisy: 10 pt, ta sama rodzina Times; Greber: „Rys.” / „Tabela”
\renewcommand{\figurename}{Rys.}
\renewcommand{\tablename}{Tabela}
\captionsetup{
  font={small,rm},
  labelfont={bf,rm},
  labelsep=period,
  justification=centering,
  skip=10pt
}

%%% 6. ALGORYTMY I LISTINGI KODU
\RequirePackage{algorithm}
\RequirePackage{algpseudocode}
\algrenewcommand\algorithmiccomment[1]{\hfill\(\triangleright\) #1}
\floatname{algorithm}{Algorytm} 

\RequirePackage{listings}
\definecolor{codebg}{RGB}{250,250,250}
\definecolor{codekw}{RGB}{0,64,160}
\definecolor{codestr}{RGB}{0,128,0}
\definecolor{codecmt}{RGB}{128,128,128}

\lstset{
  basicstyle=\ttfamily\small,
  backgroundcolor=\color{codebg},
  keywordstyle=\bfseries\color{codekw},
  stringstyle=\color{codestr},
  commentstyle=\itshape\color{codecmt},
  numbers=left,
  numberstyle=\tiny\color{gray},
  stepnumber=1,
  numbersep=8pt,
  tabsize=2,
  showstringspaces=false,
  breaklines=true,
  frame=lines,
  captionpos=b,
  % Mapowanie polskich znaków dla listings (niezbędne przy tym pakiecie)
  literate={ą}{{\k{a}}}1 {ć}{{\'c}}1 {ę}{{\k{e}}}1 {ł}{{\l{}}}1 {ń}{{\'n}}1 {ó}{{\'o}}1 {ś}{{\'s}}1 {ż}{{\.z}}1 {ź}{{\'z}}1
           {Ą}{{\k{A}}}1 {Ć}{{\'C}}1 {Ę}{{\k{E}}}1 {Ł}{{\L{}}}1 {Ń}{{\'N}}1 {Ó}{{\'O}}1 {Ś}{{\'S}}1 {Ż}{{\.Z}}1 {Ź}{{\'Z}}1
}
\renewcommand{\lstlistingname}{Listing}
\renewcommand{\lstlistlistingname}{Spis listingów}

% Styl Haskell 
\definecolor{darkgreen}{rgb}{0,0.5,0}
\lstdefinestyle{haskellStyle}{
    language=Haskell,
    commentstyle=\color{darkgreen},
    keywordstyle=\color{blue}\bfseries,
    basicstyle=\ttfamily\small,
    breaklines=true,
    stringstyle=\color{red}
}

%%% 7. HYPERREF I CLEVEREF (Musi być na końcu)
\RequirePackage[
    pdfusetitle,
    colorlinks=true,
    linkcolor=black,
    citecolor=blue,
    urlcolor=blue,
    bookmarksnumbered=true
]{hyperref}

\RequirePackage[nameinlink]{cleveref}

% Tłumaczenia dla cleveref (Greber: „rys.” / pełne „tabela”)
\crefname{figure}{rys.}{rys.}
\Crefname{figure}{Rysunek}{Rysunki}
\crefname{table}{tabela}{tabele}
\Crefname{table}{Tabela}{Tabele}
\crefname{equation}{równ.}{równ.}
\Crefname{equation}{Równanie}{Równania}
\crefname{listing}{listing}{Listingi}
\crefname{algorithm}{algorytm}{algorytmy}
\Crefname{algorithm}{Algorytm}{Algorytmy}

%%%%%%%%%%%
%% MAKRA %%
%%%%%%%%%%%

% Definicja rodzaju pracy 
\newcommand{\stRodzajPracy}{%
  \ifnum\stCzyMgr=1 MAGISTERSKA\else INŻYNIERSKA\fi
}

\newcommand{\StronaTytulowa}{
  \pagenumbering{roman}
  \normalfont % ta sama rodzina Times co reszta pracy (bez sffamily)
  \begin{titlepage}
    \newgeometry{top=2.5cm, bottom=2.5cm, left=2.5cm, right=2.5cm}

    \begin{center}
      {\fontsize{18}{22}\selectfont\bfseries Politechnika Wrocławska \par}
      \vspace{0.4cm}
      {\fontsize{14}{17}\selectfont Wydział Informatyki i Telekomunikacji \par}
    \end{center}
    \par\noindent\rule{\textwidth}{1pt}\par
    \vspace{0.5cm}

    {\normalsize
    \begin{tabular}{@{}ll}
      Kierunek: & \textbf{\stKierunek\ (\stKodKierunku)} \\
      \ifdefempty{\stSpecjalnosc}{}{Specjalność: & \textbf{\stSpecjalnosc\ (\stKodSpecjalnosci)}}
    \end{tabular}
    }

    \vspace*{\fill}

    \begin{center}
      {\fontsize{16}{19}\selectfont\bfseries PRACA DYPLOMOWA \par}
      \vspace{0.3cm}
      {\fontsize{16}{19}\selectfont\bfseries \stRodzajPracy \par}
      \vspace{2.0cm}
      {\fontsize{14}{17}\selectfont\bfseries \stTytulPL \par}
      \vspace{1cm}
      {\fontsize{14}{17}\selectfont \stAutor \par}
      \vspace{2cm}
      {\normalsize Opiekun pracy:\par \textbf{\stOpiekun} \par}
    \end{center}

    \vspace*{\fill}

    {\raggedright
      \small Słowa kluczowe: \textit{\stSlowaKluczowe}\par
      \vspace{0.5em}
      \par\noindent\rule{\textwidth}{1pt}\par
    }

    \begin{center}
      \fontsize{14}{17}\selectfont Wrocław \stRok
    \end{center}

    \restoregeometry
  \end{titlepage}
  \normalfont
  \pagestyle{empty}
}

\newenvironment{abstracts}{}{\selectlanguage{polish}\clearpage}

\newcommand{\AbstractPL}{
  \begin{center}{\WitChapterFont Streszczenie \par}\end{center}\vspace{0.5em}
}

\newcommand{\AbstractENG}{
  \selectlanguage{english}
  \begin{center}{\WitChapterFont Abstract \par}\end{center}\vspace{0.5em}
}

\newcommand{\SpisTresci}{
  \clearpage
  \addtocontents{toc}{\protect\thispagestyle{empty}} 
  \pagestyle{empty} 
  \tableofcontents
  \clearpage
  \pagestyle{empty}
}

%% WSTAWIENIE METADANYCH DO PLIKU PDF
\newcommand{\SetMetaData}{%
  \AtBeginDocument{%  
    \hypersetup{%
      pdfinfo = {%
        Title    = {\stTytulPL},
        Author   = {\stAutor},
        Subject  = {Praca dyplomowa \stRodzajPracy},  
        Keywords = {\stSlowaKluczowe}, 
        ModDate  = {\pdfcreationdate}, 
        Creator  = {pdftex}
      }
    }
  }
}
%%% KONIEC PLIKU SETTINGS.TEX